\section*{第 \ref{ch:b1:ecosystem-organization} 章——生态系统与营养结构}

\begin{solution}{exo:b1:ecosystem-organization:1}
池塘\omterm{def:b1:ecosystem-organization:ecosystem}{生态系统}即\omterm{def:b1:ecosystem-organization:ecosystem}{群落}——也就是其中全部的\omterm{def:b1:populations:population}{种群}：藻类、
眼子菜、蜗牛、昆虫幼虫、青蛙、鱼类、细菌——与
\omterm{def:b1:ecosystem-organization:ecosystem}{生境}共同构成：水体、水温、光照、溶解气体
与无机盐，以及淤泥。
\end{solution}

\begin{solution}{exo:b1:ecosystem-organization:2}
\omterm{def:b1:ecosystem-organization:niche}{栖息地}：物种生活的场所（如云杉林）。\omterm{def:b1:ecosystem-organization:niche}{生态位}：其在该处的生活方式——
即取食什么、何时取食、在树上何处、能耐受什么——
基础\omterm{def:b1:ecosystem-organization:niche}{生态位}：单独存在时可占据的范围；
实际\omterm{def:b1:ecosystem-organization:niche}{生态位}：竞争者与捕食者留给它的范围。
\end{solution}

\begin{solution}{exo:b1:ecosystem-organization:3}
$20\,800/1.7\times 10^6 = \qty{1.2}{\%}$；植食性动物
取食了净生产量的 $3400/8800 = \qty{39}{\%}$。
\end{solution}

\begin{solution}{exo:b1:ecosystem-organization:4}
\omterm{def:b1:ecosystem-organization:trophic}{生产者}、\omterm{def:b1:ecosystem-organization:trophic}{消费者}（初级、次级、三级）、\omterm{def:b1:ecosystem-organization:trophic}{分解者}、
食碎屑者；其中\omterm{def:b1:ecosystem-organization:trophic}{分解者}将死亡物质中的
无机盐归还给\omterm{def:b1:ecosystem-organization:trophic}{生产者}。
\end{solution}

\begin{solution}{exo:b1:ecosystem-organization:5}
NPP $= 800\times 17 = \qty{13.6}{MJ/m^2}$；被取食 $\qty{2.04}{MJ}$；
同化量 $\qty{1.02}{MJ}$；兔子 $\qty{51}{kJ/m^2}$，即
每公顷 \qty{510}{MJ}（\qty{30}{kg} 兔子，按 \qty{17}{kJ/g} 计）。
效率为 $51/13\,600 = \qty{0.4}{\%}$。
\end{solution}

\begin{solution}{exo:b1:ecosystem-organization:6}
生物量是存量：被取食速度与生长速度相同的浮游植物
无需积累，因此\omterm{def:b1:ecosystem-organization:trophic}{生产者}的现存量可以
少于以其为食的寿命更长的动物的现存量。能量则是
流量：每一\omterm{def:b1:ecosystem-organization:trophic}{营养级}从下一级获得的全部能量
都须有一部分以热量形式散失，因此能量流只能逐级减少。
\end{solution}

\begin{solution}{exo:b1:ecosystem-organization:7}
$p = 0.5, 0.3, 0.15, 0.05$: $H = -(0.5\ln 0.5 + 0.3\ln 0.3 + 0.15\ln
0.15 + 0.05\ln 0.05) = 0.347 + 0.361 + 0.285 + 0.150 = 1.14$;
均匀度为 $1.14/\ln 4 = 0.82$。四者各为 25：$H = \ln 4 = 1.39$，
均匀度为 1。
\end{solution}

\begin{solution}{exo:b1:ecosystem-organization:8}
净值 $= \qty{3}{mg}$ 的 $\mathrm{O_2}$，每升每天；呼吸作用
\qty{1}{mg}；总值 $3 + 1 = \qty{4}{mg}$。以碳计，总值为
\qty{1.5}{mg}，净值为 \qty{1.1}{mg}，均为每升每天的 C。
\end{solution}

\begin{solution}{exo:b1:ecosystem-organization:9}
每一环节保留十分之一：经过十个环节后，只剩下初级生产量的
十亿分之一，在任何面积上都不足以养活一只动物。
生产力低（如沙漠）或面积小（如岛屿）时，
十分之一的十分之一在更少的环节后即告耗尽；森林的
高生产力可多支持一两个环节。
\end{solution}

\begin{solution}{exo:b1:ecosystem-organization:10}
谷物：\qty{1}{t} 需要六分之一公顷的土地。牛肉：\qty{1}{t} 牛肉所含的
能量需 \qty{10}{t} 谷物、占地 \num{1.7} 公顷，
即十倍于前者的土地。以第二\omterm{def:b1:ecosystem-organization:trophic}{营养级}为食，需多付出一个\omterm{def:b1:ecosystem-organization:trophic}{营养级}的能量；
一颗星球以谷物直接养人，所能养活的人数是
吃谷物喂养的肉类时的十倍，而人类按各种饮食平均
所处的\omterm{def:b1:ecosystem-organization:trophic}{营养级}，是\omterm{thm:b1:populations:logistic}{环境容纳量}中的一个重要项。
\end{solution}

\begin{solution}{exo:b1:ecosystem-organization:11}
以每平方米计，海洋是一片荒漠，但它覆盖了地球的三分之二，
因此低速率乘以巨大面积，便得到
总量的一半。其生产力受养分限制，养分
沉入光照的表层；洋流将养分带回上层的地方（沿岸上升流、
大陆架），生产力高十倍；在这些地方，仅占
海洋的百分之几，鱼类与渔业便高度集中。
\end{solution}

\begin{solution}{exo:b1:ecosystem-organization:12}
能量以光的形式进入，流经若干\omterm{def:b1:ecosystem-organization:trophic}{营养级}的生物，
最终全部以热量的形式散失。从这个意义上说，\omterm{def:b1:ecosystem-organization:ecosystem}{生态系统}确实把
阳光转化为热量；\omterm{def:b1:ecosystem-organization:trophic}{分解者}呼吸\omterm{def:b1:ecosystem-organization:trophic}{消费者}未利用的一切，
是大部分能量的去向。但物质并未被消耗：
碳、氮和磷循环往复，同样由\omterm{def:b1:ecosystem-organization:trophic}{分解者}
归还给\omterm{def:b1:ecosystem-organization:trophic}{生产者}，因此系统得以持续；而
\omterm{def:b1:organism-environment:organism}{生物体}并非副产品，而是能量流动借以\omterm{def:b1:body-plans-tissues:tissue}{组织}的结构
——\omterm{def:b1:ecosystem-organization:niche}{生态位}、\omterm{def:b1:ecosystem-organization:foodweb}{食物网}与金字塔，正是能量在流动中
所构建的东西。这句话讲对了热力学，
却漏掉了生物学。
\end{solution}

\begin{solution}{pb:b1:ecosystem-organization:1}
\textbf{1.} $20\,800 - 12\,000 = \qty{8800}{kcal}$，每平方米每
年。
\textbf{2.} \qty{1.2}{\%}。损失包括：未被吸收的光（光谱的一半、
反射和水体吸收）；叶片在全日照或遮阴下已饱和，
光呼吸，以及植物自身的呼吸（已计入
总值与净值之间）。
\textbf{3.} $12\,000/20\,800 = \qty{58}{\%}$。
\textbf{4.} 被取食 \qty{3400}{}；
直接流向\omterm{def:b1:ecosystem-organization:trophic}{分解者}：$8800 - 3400 = \qty{5400}{kcal}$。
\textbf{5.} $8800/10 = \qty{880}{g/m^2}$：与温带森林或
肥沃草原相当。
\textbf{6.} $1.7\times 10^6 - 20\,800 = \qty{1.68e6}{kcal}$，占 \qty{98.8}{\%}
的能量被反射、透射，或被水与岩石吸收后
以热量形式再辐射。
\textbf{7.} 生产量 $1500 - 1100 = \qty{400}{kcal}$；同化
效率为 $1500/3400 = \qty{44}{\%}$。
\textbf{8.} $400/1500 = \qty{27}{\%}$。
\textbf{9.} $400/8800 = \qty{4.5}{\%}$（若按
被取食量计，即 \qty{3400}{}，则为 \qty{16}{\%}）。
\textbf{10.} 肉食动物的生产量 $380 - 300 = \qty{80}{kcal}$；
效率为 $80/400 = \qty{20}{\%}$（它们取食了所产的
\qty{400}{} 中的 \qty{380}{}）。
\textbf{11.} $6/80 = \qty{7.5}{\%}$。
\textbf{12.} \qty{6}{kcal}；$6/1.7\times 10^6 = \num{3.5e-6}$，约为
百万分之几。
\textbf{13.} 肉类易于消化，且成分与取食者接近；
植物物质为纤维状，大多难以消化，因此食草动物
以粪便形式损失的摄入量更多。
\textbf{14.} 未被取食的净生产量 \qty{5400}{}；食草动物粪便
$3400 - 1500 = \qty{1900}{}$；食草动物未被取食的生产量
$400 - 380 = \qty{20}{}$；肉食动物粪便（未给出同化量；
按被取食量减去呼吸量与生产量 $= 0$ 计，即全部被同化）
以及未被取食的肉食动物生产量 $80 - 20 = \qty{60}{}$；顶级肉食动物
的生产量 \qty{6}{}：合计 \qty{7386}{kcal}。
\textbf{15.} $12\,000 + 1100 + 300 + 14 + 7386 = \qty{20800}{kcal}$。
\textbf{16.} 与总初级生产量相等：泉处于
\omterm{def:b1:organism-environment:homeostasis}{稳态}，年际间不储存生物量。
\textbf{17.} $7386/20\,800 = \qty{36}{\%}$（其余大部分为
\omterm{def:b1:ecosystem-organization:trophic}{生产者}自身的呼吸）。
\textbf{18.} \omterm{def:b1:ecosystem-organization:trophic}{生产者} \qty{8800}{}，食草动物 \qty{400}{}，肉食动物
\qty{80}{}，顶级肉食动物 \qty{6}{}；效率分别为 \qty{4.5}{\%}，
\qty{20}{\%}、\qty{7.5}{\%}。
\textbf{19.} \qty{8000}{}、\qty{400}{}、\qty{100}{}、\qty{15}{kcal/m^2}：
呈正金字塔形，每一\omterm{def:b1:ecosystem-organization:trophic}{营养级}为下一级的十分之一或更少。
\textbf{20.} \omterm{def:b1:ecosystem-organization:trophic}{生产者} $8000/8800 = 0.9$ 年；食草动物 $400/400 =
1$ 年；肉食动物 $100/80 = 1.25$ 年；顶级肉食动物 $15/6 =
2.5$ 年。植物更替最快：\omterm{def:b1:body-plans-tissues:tissue}{组织}寿命短，被
持续取食；顶级肉食动物是寿命长的动物，其存量
更新缓慢。
\textbf{21.} 一只 \qty{2}{kg} 的鹭干重约 \qty{500}{g}，
\qty{5000}{kcal}，每年需要约相当的生产量，以
替换自身及其幼体：$5000/6 \approx \qty{800}{m^2}$ 的
泉域——实际远不止如此，因其还需呼吸：按每天
\qty{100}{kcal}，全年摄入 \qty{36500}{kcal}，即
\qty{6000}{m^2} 的下一\omterm{def:b1:ecosystem-organization:trophic}{营养级}（肉食动物
生产量 \qty{80}{}），这便是鹭领地的样子。
\textbf{22.} 多出的净生产量未被取食，流向
\omterm{def:b1:ecosystem-organization:trophic}{分解者}，其呼吸作用加倍；夜间没有
光合作用，它们的耗氧需求可耗尽水中的氧气，
导致鱼类死亡——即富营养化。
\textbf{23.} 食草动物的生产量降至 \qty{200}{}：肉食动物
获得的食物减半，其生产量降至约 \qty{40}{}，顶级肉食动物
的生产量降至约 3，鹭的数量也随之减少；鳗草
被取食得少而更加茂密，更多的部分未被取食而流向\omterm{def:b1:ecosystem-organization:trophic}{分解者}。
\textbf{24.} 顶级肉食动物的能量很少，但其影响
体现在肉食动物的数量上：这一数量控制食草动物，
食草动物又控制草：移除它们会使肉食动物增多，
食草动物减少、草生长旺盛——这是\omterm{def:b1:ecosystem-organization:foodweb}{食物网}中的
级联（见\cref{ch:b1:species-interactions}）；移除几千卡的
食草动物，结果只改变了几千卡能量。
\textbf{25.} \omterm{def:b1:ecosystem-organization:trophic}{生产者}到食草动物 \qty{4.5}{\%}（占 NPP；
占被取食量的 \qty{16}{\%}）；食草动物到肉食动物 \qty{20}{\%}，
肉食动物到顶级肉食动物 \qty{7.5}{\%}；阳光到总初级生产量
\qty{1.2}{\%}；仅百万分之几的阳光最终流入顶级
肉食动物。
\end{solution}
